\section{Stiefel incidence and conditional block information}
\label{sec:incidence}

This comparison bounds conditional information directly by the relative
entropy of two laws on tuples of signals.  Both laws make the signals share
an exact projection.  The first uses the actual observation block; the
second obtains its shared projection from an orthonormal frame inside the
coordinates of the independent side projection.  Given the separate side
projections of all the signals, the second law leaves them conditionally
independent with their original fiber laws.  This product structure makes
the divergence of the two laws an upper bound for the information in the
message, multiplied by the number of signals.

The geometric work is to express both laws relative to one incidence
measure.  We first construct the laws and prove their information
comparison.  We then compute their exact Jacobians, transfer density
exceptions before division, and estimate the resulting likelihood ratio.

Write \(\sigma=\sigma_d\), and let \(v_q\) denote the volume of the
unit ball in \(\mathbb R^q\).  For the conditional chain rule on standard
Borel spaces, see
Austin~\cite[Sections~3.3--3.4]{Austin2020}.  We use, for all sufficiently
large \(d\),
\begin{equation}\label{inc:dimensions}
 n=k=\lfloor d/16\rfloor,\qquad
 \ell=k-1,\qquad p=\lfloor d/4\rfloor .
\end{equation}
Thus \(\ell\ge2\), while \(p-n-\ell\) and \(d-n-\ell\) are at least an
absolute positive multiple of \(d\).

\begin{proposition}[Information in one block]\label{inc:block}
Let \(\mu=f\sigma\) be a probability measure with a Borel density
\(0\le f\le F<\infty\), where \(F\ge1\).  Draw \(S\sim\mu\) and standard
Gaussian matrices \(A\in\mathbb R^{n\times d}\) and
\(B\in\mathbb R^{p\times d}\), all independently.  For an integer
\(w\ge1\), a message \(J\) takes values in a fixed \(w\)-point alphabet
(unused symbols are allowed) and is drawn by an arbitrary measurable
probability kernel from
\((A,AS)\); conditional on \((A,AS)\), its randomness is independent of
\((S,B)\).  There is an absolute constant \(C\) such that
\begin{equation}\label{inc:block-bound}
 I_\mu(S;J\mid B,BS)
 \le \frac{\log w}{k}+Cd+C\log(1+\log F).
\end{equation}
\end{proposition}

The dependence on \(F\) is doubly logarithmic.  This permits us later to
condition on an entire finite history: a history of probability \(a\)
changes the uniform prior to a density bounded by \(1/a\).

\subsection{Two probability laws on compatible signals}
\label{inc:laws}

We use the spherical projection versions from
\eqref{eq:finite-tilt}, with prior \(\mu=f\sigma\).  Recall them locally:
for a full-row-rank \(q\times d\) matrix \(L\), with
\(1\le q\le d-2\), and an interior label \(t\), let
\(\sigma_{L,t}\) be uniform probability on the spherical fiber \(Ls=t\),
and let \(p_L^0\) be the uniform-sphere projection density in
\eqref{eq:spherical-projection-density}.  Then
\begin{equation}\label{inc:recalled-fibers}
 p_L(t)=p_L^0(t)\int f\,d\sigma_{L,t},\qquad
 \mu_{L,t}=\frac{f\,\sigma_{L,t}}{\int f\,d\sigma_{L,t}}
 \quad\text{when }p_L(t)>0.
\end{equation}
Set \(p_L=0\) outside the interior projection ellipsoid and fix one
measurable probability kernel on unused fibers.  These jointly measurable
versions specify all the conditional draws below.  Their geometric
finite-measure form will enter only when we compute the two likelihoods.

Write \(\phi_j(A,y)\) for the probabilities of the message kernel.
For a tuple \(\mathbf s=(s_1,\ldots,s_k)\), put
\[
 U=[s_2-s_1,\ldots,s_k-s_1]\in\mathbb R^{d\times\ell}.
\]
The notation \(U^\perp\) means the orthogonal complement of the column
space of \(U\), and \(P_{U^\perp}\) is its orthogonal projector.
Let \(\mathcal V_{n,h}\) be the manifold of \(n\times h\) matrices with
orthonormal rows, with the metric induced by the Euclidean matrix norm.
Write \(\upsilon_{n,h}\) for normalized Riemannian volume on it.
This is the Euclidean Stiefel metric discussed by Edelman, Arias, and
Smith~\cite[Section~2.2.1]{EdelmanAriasSmith1998}; we compute its volume
factor below.

We use an auxiliary \(C\in\mathcal V_{n,p}\) and set \(D=CB\).
For a full-rank \(p\times\ell\) matrix \(W\), write \(\upsilon_W\) for
uniform frames whose rows lie in \((\operatorname{col}W)^\perp\).
Write \(\gamma_{r,d}\) for standard Gaussian measure on
\(\mathbb R^{r\times d}\), and \(\gamma_U\) for the law of \(n\)
independent \(N(0,P_{U^\perp})\) rows.  All these kernels are measurable
on their rank domains.  For example, projecting a Gaussian matrix onto
\((\operatorname{col}W)^\perp\) and taking its orthonormal row polar
factor gives \(\upsilon_W\); orthogonal invariance identifies this with
normalized induced volume.  We may assign arbitrary values on rank
exceptions.

Define \(Q\) by the following experiment.
\begin{enumerate}
\item Draw \(A\sim\gamma_{n,d}\) and \(s_1\sim\mu\) independently.
      Put \(y=As_1\).  Draw \(J\) from \(\phi(A,y)\) and draw
      \(s_2,\ldots,s_k\) independently from \(\mu_{A,y}\), independently
      of \(J\) conditional on \((A,y)\).
\item Draw \(B\sim\gamma_{p,d}\) independently.  Given \((\mathbf s,B)\),
      draw \(C\sim\upsilon_{BU}\).
\end{enumerate}
Then \(AU=0\).  Conditional on \((A,y)\), all \(k\) signals are
independent with law \(\mu_{A,y}\), which is absolutely continuous on
a positive-radius sphere of dimension \(d-n-1\).  Each new point
therefore avoids the affine span of its predecessors; hence \(U\) has
rank \(\ell\) almost surely.  The independent Gaussian \(B\) makes
\(BU\) have rank \(\ell\) almost surely, so \(\upsilon_{BU}\) is
available because \(p-\ell\ge n\).  Each marginal
\begin{equation}\label{inc:q-marginals}
 \mathcal L_Q(s_i,J,B,Bs_i)=\mathcal L_\mu(S,J,B,BS)
 \qquad(1\le i\le k)
\end{equation}
is the law in Proposition~\ref{inc:block}.

The matrix \(D=CB\) has a useful conditional law under \(Q\):
\begin{equation}\label{inc:conditional-D}
 \mathcal L_Q(D\mid \mathbf s,A,y,J)
 =N(0,P_{U^\perp})^{\otimes n}.
\end{equation}
To verify this, split each row of \(B\) between
\(\operatorname{col}U\) and \(U^\perp\).  The second component is
independent of \(BU\).  Given the tuple, the kernel for \(C\) depends
on \(B\) only through \(BU\); its orthonormal rows send the independent
Gaussian second component to \(n\) independent standard Gaussian rows
in \(U^\perp\).  The first component is annihilated by \(C\).

Define \(R\) in the other order.  First draw independent
\(B\sim\gamma_{p,d}\) and \(C\sim\upsilon_{n,p}\); then \(D=CB\) is
standard Gaussian.  Draw \(s_1\sim\mu\) independently and, conditional
on \((D,Ds_1)\), draw \(s_2,\ldots,s_k\) independently from
\(\mu_{D,Ds_1}\).  Draw \(J\) uniformly from its \(w\) possible values,
independently, and finally draw \(A\sim\gamma_U\).
Again \(U\) has rank \(\ell\).  Conditional on \((C,D)\), the other
\(p-n\) rows in an orthogonal row rotation of \(B\) are independent
standard Gaussian rows and were not used to draw the tuple.  Since
\(p-n\ge\ell\), they make \(BU\) have rank \(\ell\) almost surely.
In these experiments \(A\) and \(D\) also have row rank \(n\) almost
surely: each is either standard Gaussian or has conditional Gaussian
support \(U^\perp\), of dimension \(d-\ell\ge n\).

\begin{lemma}[Product conditioning]\label{inc:product}
Put \(Z_i=Bs_i\).  After \(A,C\) are discarded under \(R\),
\begin{equation}\label{inc:product-law}
 \mathcal L_R(\mathbf s\mid B,Z_1,\ldots,Z_k,J)
 =\bigotimes_{i=1}^k\mu_{B,Z_i}.
\end{equation}
Consequently,
\begin{equation}\label{inc:entropy-lower}
 \operatorname{KL}(Q\|R)\ge kI_\mu(S;J\mid B,BS).
\end{equation}
\end{lemma}

\begin{proof}
Retain \(B,C\) temporarily and put \(t=Ds_1\).  Conditional on
\((B,C,t)\), the signals are independent with law \(\mu_{D,t}\).
Disintegrating this law further according to \(z=Bs\) gives
\(\mu_{B,z}\): the earlier equation \(Ds=t\) is just \(Cz=t\),
already determined by \(z\).  Formally, disintegrate the law
of \(z=Bs\) according to \(Cz\), and compose its conditional kernel
with \(\mu_{B,z}\); the resulting kernel is a conditional law given
\(Ds=t\).  These identities hold at the needed values because each
individual signal has marginal law \(\mu\) given \((B,C)\).
Conditioning the independent signals on their separate \(Z_i\)'s
therefore yields the product in \eqref{inc:product-law}.  That product
does not depend on \(C,t\), so discarding them leaves it unchanged.
The independent uniform \(J\) and the later kernel for \(A\) do not
change the marginal under consideration.

Data processing and the conditional chain rule give
\[
 \operatorname{KL}(Q\|R)\ge
 \mathbb E_Q\operatorname{KL}\!\left(
 Q(\mathbf s\mid B,Z_1,\ldots,Z_k,J)
 \,\middle\|\,\bigotimes_{i=1}^k\mu_{B,Z_i}\right).
\]
This is valid as an extended inequality; absolute continuity will be
established below.  Relative entropy to a product is at least the sum
of the relative entropies of the marginals.  For each such summand,
convexity of relative entropy permits us, after averaging, to discard
the other \(Z_j\)'s from its first conditional law.  The result is at
least
\[
 \mathbb E_Q\operatorname{KL}\bigl(
 Q(s_i\mid B,Z_i,J)\,\|\,\mu_{B,Z_i}\bigr)
 =I_\mu(S;J\mid B,BS)
\]
by \eqref{inc:q-marginals}.  Sum over \(i\).
\end{proof}

The product conditional law has reduced the block information bound to
one task: bound \(\operatorname{KL}(Q\|R)\) by \(\log w\) plus the
geometric and prior-density costs.  We next put both laws on a common
incidence measure.  Its exact densities will make that divergence
computable, after exceptional sets have been removed before division.

\subsection{The common measure and its density versions}
\label{inc:common-measure}\label{inc:fibers}

To compare the two laws, we express their conditional probabilities as
normalized finite measures on the same incidence set.  The underlying
geometric formula is the regular-submersion case of coarea; see
Federer~\cite[Definition~2.10 and Theorem~3.1, pp.~423 and 426--427]{Federer1959}.
The following calculation specifies the induced fiber volume and the
probability normalization used here.

\begin{lemma}[Spherical projection measure]\label{inc:projection}
Suppose \(L\in\mathbb R^{q\times d}\) has full row rank and
\(1\le q\le d-2\).  In the interior
\[
 E_L=\{t:t^{\mathsf T}(LL^{\mathsf T})^{-1}t<1\}
\]
define a finite measure on \(\{s\in\mathbb S^{d-1}:Ls=t\}\) by
\begin{equation}\label{inc:fiber-measure}
 \lambda_{L,t}(ds)=
 \frac{f(s)}{|\mathbb S^{d-1}|\,j_L(s)}
 \,d\operatorname{vol}_{d-1-q}(s),\qquad
 j_L(s)=\det\!\bigl(L(I-ss^{\mathsf T})L^{\mathsf T}\bigr)^{1/2}.
\end{equation}
Set \(\lambda_{L,t}=0\) outside \(E_L\).  Its mass and normalization
are exactly the versions recalled in \eqref{inc:recalled-fibers}:
\[
 p_L(t)=\lambda_{L,t}(\mathbb S^{d-1}),\qquad
 \mu_{L,t}=\lambda_{L,t}/p_L(t)\quad\text{where }p_L(t)>0.
\]
The finite-measure kernel is jointly measurable in \((L,t)\).  Moreover,
\begin{equation}\label{inc:ellipsoid-bound}
 V_L\|p_L\|_\infty\le Fd^d,\qquad
 V_L=|E_L|=v_q\det(LL^{\mathsf T})^{1/2}.
\end{equation}
\end{lemma}

\begin{proof}
The tangent space at \(s\) is \(s^\perp\).  The Gram matrix of the
differential \(L:s^\perp\to\mathbb R^q\) is
\(L(I-ss^{\mathsf T})L^{\mathsf T}\), proving the formula for its normal
Jacobian.  With
\(P_L=L^{\mathsf T}(LL^{\mathsf T})^{-1}L\), the determinant lemma also gives
\[
 j_L(s)=\det(LL^{\mathsf T})^{1/2}\sqrt{1-\|P_Ls\|^2}.
\]
The differential is onto except when \(s\) belongs to the row space of
\(L\); that sphere has zero \((d-1)\)-dimensional surface measure.
Applying change of variables on coordinate patches of its complement,
and then exhausting the complement, gives
\[
 \int h(s)\,d\mu(s)
 =\int_{\mathbb R^q}\int h(s)\,\lambda_{L,t}(ds)\,dt
 \quad(h\ge0).
\]
To identify these measures with the versions already used, put
\(s_0=L^{\mathsf T}(LL^{\mathsf T})^{-1}t\) and
\(r=\sqrt{1-\|s_0\|^2}\).  An interior fiber is
\(s_0+r\mathbb S(\ker L)\), where \(\mathbb S(\ker L)\) is the unit
sphere of \(\ker L\), and \(j_L\) is constant there.  Its induced
volume is \(|\mathbb S^{d-q-1}|r^{d-q-1}\).  Dividing this volume by
\(|\mathbb S^{d-1}|j_L\) gives precisely the uniform projection density
\(p_L^0(t)\) in \eqref{eq:spherical-projection-density}.  Consequently,
as finite measures on every interior fiber,
\[
 \lambda_{L,t}=p_L^0(t)f\,\sigma_{L,t}.
\]
Taking mass and then normalizing gives \eqref{inc:recalled-fibers}
pointwise on its positive set.  A uniform
direction in \(\ker L\) can be generated by projecting a standard
Gaussian vector onto that subspace and normalizing.  The projection
matrix is a measurable function of full-rank \(L\); integrating the
bounded Borel function \(f(s_0+r\theta)\) against this kernel therefore
gives measurable \(\lambda_{L,t}\), \(p_L(t)\), and \(\mu_{L,t}\),
with the fixed convention on unused fibers from \eqref{inc:recalled-fibers}.

For orthonormal rows and \(f=1\), the displayed fiber formula gives a
density proportional to
\((1-\|t\|^2)^{(d-q-2)/2}\) on the unit ball.  Its normalizing integral
is at least \(\tfrac12v_qd^{-q/2}\): restrict to the ball of radius
\(d^{-1/2}\) and use
\((1-1/d)^{(d-q-2)/2}\ge1/2\).
The maximum of the density times \(v_q\) is consequently at most
\(2d^{q/2}\le d^d\).  The bound \(f\le F\) multiplies this estimate by
at most \(F\).  Finally, an invertible change from orthonormal
row coordinates to \(L\) divides the density by
\(\det(LL^{\mathsf T})^{1/2}\) and multiplies the ellipsoid volume by
the same factor.  This proves \eqref{inc:ellipsoid-bound}; in particular
our chosen \(p_L\) is finite everywhere.
\end{proof}


\begin{lemma}[Interchange on a regular incidence set]\label{inc:swap}
Let \(X,Z\) be second-countable Riemannian manifolds, and let
\(G:X\times Z\to\mathbb R^a\) be smooth.  On a relatively open subset
\(\Omega\) of \(G^{-1}(0)\), suppose both partial differentials
\(d_xG\) and \(d_zG\) are onto.  Write \(J_xG,J_zG\) for their normal
Jacobians.  For every nonnegative measurable \(h\) on \(\Omega\),
\begin{align}
 &\int_Z\int_{\{x:(x,z)\in\Omega\}}
       \frac{h(x,z)}{J_xG(x,z)}
       \,d\operatorname{vol}_{\dim X-a}(x)\,d\operatorname{vol}_Z(z)
       \notag\\
 &\qquad=
 \int_X\int_{\{z:(x,z)\in\Omega\}}
       \frac{h(x,z)}{J_zG(x,z)}
       \,d\operatorname{vol}_{\dim Z-a}(z)\,d\operatorname{vol}_X(x).
 \label{inc:swap-identity}
\end{align}
The fiber volumes are induced volumes, and the equality holds for
nonnegative integrals even when their common value is infinite.
\end{lemma}

\begin{proof}
Near a point of \(\Omega\), choose coordinates
\(x=(x',x'')\), \(z=(z',z'')\) so that \(x''\) and \(z''\) each have
\(a\) coordinates and both \(G_{x''}\) and \(G_{z''}\) are invertible.
Let \(u_X(x)\,dx\), \(u_Z(z)\,dz\) be the volume forms in these
coordinates.  Changing variables from \((x',x'')\) to \((x',G)\)
shows that the induced \(x\)-fiber volume divided by \(J_xG\) is
\[
 \frac{u_X(x)}{|\det G_{x''}|}\,dx'.
\]
Indeed, the tangent-volume factor and the normal-volume factor of this
change of variables are respectively the induced fiber volume and the
reciprocal normal Jacobian.  Thus the left integral has local density
\(u_Xu_Z/|\det G_{x''}|\) in \((x',z',z'')\); the right integral has
density \(u_Xu_Z/|\det G_{z''}|\) in \((x',z',x'')\).
At fixed \((x',z')\), implicit differentiation gives
\[
 \frac{dx''}{dz''}=-G_{x''}^{-1}G_{z''},\qquad
 \left|\det\frac{dx''}{dz''}\right|
 =\frac{|\det G_{z''}|}{|\det G_{x''}|}.
\]
The two local integrals agree.  Split \(\Omega\) into measurable pieces
subordinate to a countable coordinate cover and sum using Tonelli's
theorem.  This also proves the assertion for infinite integrals.
\end{proof}

Let \(\mathsf H\) be the probability law that draws
\(\mathbf s\sim\mu^{\otimes k}\), then independent
\(B\sim\gamma_{p,d}\), and, conditional on these, independently draws
\(A\sim\gamma_U\) and \(C\sim\upsilon_{BU}\).  Include an independent
uniform \(J\).  Independent points from \(\mu\) give
\(\operatorname{rank}U=\ell\) almost surely.  The same Gaussian
decomposition as in \eqref{inc:conditional-D} then gives full ranks for
\(BU,A,D\) under \(\mathsf H\).

\begin{lemma}[Exact undivided identities and probability densities]
\label{inc:densities}
Define the following marginal probability measures of \(\mathsf H\):
\[
 \begin{split}
 \mathsf H_A(d\mathbf s,dA)
   &=\mu^{\otimes k}(d\mathbf s)\,\gamma_U(dA),\\
 \mathsf H_C(d\mathbf s,dB,dC)
   &=\mu^{\otimes k}(d\mathbf s)\,\gamma_{p,d}(dB)\,\upsilon_{BU}(dC).
 \end{split}
\]
Using the unnormalized measures of \eqref{inc:fiber-measure}, define
nonnegative measures on the corresponding exact incidence sets by
\begin{align}
 \mathsf M_A(d\mathbf s,dA)
  &=\mu(ds_1)\gamma_{n,d}(dA)
       \prod_{i=2}^k\lambda_{A,As_1}(ds_i),\label{inc:MA}\\
 \mathsf M_C(d\mathbf s,dB,dC)
  &=\mu(ds_1)\gamma_{p,d}(dB)\upsilon_{n,p}(dC)
       \prod_{i=2}^k\lambda_{D,Ds_1}(ds_i),\qquad D=CB.\label{inc:MC}
\end{align}
Values at row-rank failures may be set to zero.  The weights on the
right below are assigned arbitrary finite values on their Gram-rank
null sets.  These measures satisfy the exact identities
\begin{align}
 \mathsf M_A
  &=(2\pi)^{-n\ell/2}\det(U^{\mathsf T}U)^{-n/2}\,\mathsf H_A,
       \label{inc:undivided-A}\\
 \mathsf M_C
  &=\kappa\det(U^{\mathsf T}B^{\mathsf T}BU)^{-n/2}\,\mathsf H_C,
 \qquad
 \kappa=\frac{|\mathcal V_{n,p-\ell}|}{|\mathcal V_{n,p}|}.
       \label{inc:undivided-C}
\end{align}
These are equalities of nonnegative measures.

For \(L=A,D\), the values \(p_L(Ls_1)\) are positive and finite
\(\mathsf H\)-almost surely and equal the shrinking-ball limits
\[
 p_L(Ls_1)=\lim_{j\to\infty}
 \frac{\mu\{z:\|Lz-Ls_1\|<2^{-j}\}}{v_n2^{-jn}}.
\]
On this full-measure set, the densities of \(Q\) and \(R\) with respect
to \(\mathsf H\) are
\begin{align}
 q_*&=w\phi_J(A,As_1)(2\pi)^{-n\ell/2}
       \det(U^{\mathsf T}U)^{-n/2}p_A(As_1)^{-\ell},
       \label{inc:q-density}\\
 r_*&=\kappa\det(U^{\mathsf T}B^{\mathsf T}BU)^{-n/2}
       p_D(Ds_1)^{-\ell}.\label{inc:r-density}
\end{align}
In particular \(Q\ll R\).
\end{lemma}

\begin{proof}
We first establish the identities without dividing by any projection
density.  In \(\mathsf M_A\), integrating the tuple variables gives the
exact \((s_1,A)\)-marginal
\begin{equation}\label{inc:MA-marginal}
 \mu(ds_1)\gamma_{n,d}(dA)\,p_A(As_1)^\ell.
\end{equation}
For fixed full-rank \(A\), this weight is finite by
\eqref{inc:ellipsoid-bound}.  Where it is positive,
\(\lambda_{A,As_1}^{\otimes\ell}
=p_A(As_1)^\ell\mu_{A,As_1}^{\otimes\ell}\).
The rank argument for the \(Q\) experiment therefore shows that
\(\operatorname{rank}U=\ell\) outside an \(\mathsf M_A\)-null set.
At zero weight the fiber product has zero mass.  The same reasoning in
\(\mathsf M_C\) gives the exact \((s_1,B,C)\)-marginal
\begin{equation}\label{inc:MC-marginal}
 \mu(ds_1)\gamma_{p,d}(dB)\upsilon_{n,p}(dC)\,p_D(Ds_1)^\ell,
\end{equation}
and the rank argument for \(R\) removes the loci where \(U\) or \(BU\)
has rank less than \(\ell\).  Multiplication by the finite weight in
\eqref{inc:MC-marginal} preserves these native null sets.

On either incidence set, the spherical partial differentials are
regular on the remaining locus.  Indeed, if some \(s_i\) belonged to
the row space of \(A\), then its row-space projection would have norm
one.  The equations \(As_j=As_i\) would give every \(s_j\) that same
row-space projection, and their unit norms would force \(s_j=s_i\)
for all \(j\), contradicting \(\operatorname{rank}U=\ell\).
The same argument applies to \(D\).  The Gaussian rank assertions
above also remove row-rank failures.  Thus removing the complements of
the full-rank regular loci changes neither \(\mathsf M_A,\mathsf M_C\)
nor the corresponding \(\mathsf H_A,\mathsf H_C\).

Fix \(s_1\) and apply Lemma~\ref{inc:swap} to
\[
 x=(s_2,\ldots,s_k),\qquad z=A,\qquad G(x,z)=AU.
\]
For fixed \(A\), the target has \(\ell\) independent \(n\)-coordinate
blocks, so the partial normal Jacobian is
\(\prod_{i=2}^k j_A(s_i)\).  These are precisely the denominators in
\(\prod_{i=2}^k\lambda_{A,As_1}(ds_i)\).
For fixed tuple, the differential is \(\Delta A\mapsto\Delta A\,U\).
It consists of \(n\) row maps with Gram matrix \(U^{\mathsf T}U\);
its normal Jacobian is \(\det(U^{\mathsf T}U)^{n/2}\).
The \(A\)-fiber has rows in \(U^\perp\).  Restricting the ambient
standard Gaussian density to this subspace and then expressing it
relative to the normalized Gaussian \(\gamma_U\) contributes
\((2\pi)^{-n\ell/2}\).  This proves \eqref{inc:undivided-A} on the
regular locus, and hence exactly after the null loci just proved are
removed.  In particular, integrating over a restricted regular tuple
fiber gives at most \(p_A(As_1)^\ell\), and gives equality for native
almost every \((s_1,A)\).

For the second identity, fix \((s_1,B)\) and use
\(G(x,C)=CBU\).  Put \(W=BU\).  The tuple partial normal Jacobian is
\(\prod_{i=2}^k j_D(s_i)\).  At \(CW=0\),
\[
 T_C\mathcal V_{n,p}
 =\{\Delta:\Delta C^{\mathsf T}+C\Delta^{\mathsf T}=0\}.
\]
Variations whose rows belong to \(\operatorname{col}W\) lie in this
tangent space, since \(C W=0\).  They form the orthogonal complement
of the tangent space to the constrained frame fiber.  On this normal
space the differential is \(\Delta\mapsto\Delta W\), with normal
Jacobian
\(\det(W^{\mathsf T}W)^{n/2}\).  There is no additional metric factor
in these normal directions.  The constrained fiber, with its induced
metric, is isometric to \(\mathcal V_{n,p-\ell}\).  Replacing its
unnormalized volume by \(\upsilon_W\) leaves the volume ratio \(\kappa\).
Lemma~\ref{inc:swap} now proves \eqref{inc:undivided-C}, again exactly
after removal of native rank-null sets.  The corresponding restricted
tuple mass is at most \(p_D(Ds_1)^\ell\) and equals it for native almost
every \((s_1,B,C)\).

We next transfer the density exceptions using these undivided measures.
For a full-rank \(n\times d\) matrix \(L\), define
\[
 b_j(L,t)=
 \frac{\mu\{z:\|Lz-t\|<2^{-j}\}}{v_n2^{-jn}},
\]
and let \(\mathcal E(L,s)\) be the condition that either
\(p_L(Ls)\notin(0,\infty)\), or the limit of \(b_j(L,Ls)\) does not
exist, or that limit differs from \(p_L(Ls)\).  This is a measurable
condition.  Set
\[
 \mathcal E_A=\{(\mathbf s,A):\mathcal E(A,s_1)\},\qquad
 \mathcal E_D=\{(\mathbf s,B,C):\mathcal E(CB,s_1)\}.
\]
For every fixed full-rank \(L\), Lebesgue differentiation and the
density formula imply \(\mu\{s:\mathcal E(L,s)\}=0\).
The exact native marginals \eqref{inc:MA-marginal} and
\eqref{inc:MC-marginal}, whose density weights are finite for each
fixed native projection, therefore give
\[
 \mathsf M_A(\mathcal E_A)=0,\qquad
 \mathsf M_C(\mathcal E_D)=0.
\]
On the right sides of \eqref{inc:undivided-A} and
\eqref{inc:undivided-C}, the multiplying Gram weights are strictly
positive and finite almost surely.  If \(h\) denotes either weight,
then on \(\{1/m\le h\le m\}\) the equality
\(\int_{\mathcal E}h\,d\mathsf H_\bullet=0\) gives
\(\mathsf H_\bullet(\mathcal E\cap\{1/m\le h\le m\})=0\).
Letting \(m\to\infty\) proves that \(\mathcal E_A\) is
\(\mathsf H_A\)-null and \(\mathcal E_D\) is \(\mathsf H_C\)-null.
This argument is valid even if an undivided measure has infinite total
mass.  Adding the unused probability kernels gives both assertions
under \(\mathsf H\).

Only now replace each unnormalized fiber measure by its normalized
conditional law.  The \(Q\) experiment divides
\eqref{inc:undivided-A} by \(p_A(As_1)^\ell\), then attaches the same
kernels for \(B,C\) and the message density \(w\phi_J(A,As_1)\) relative
to uniform \(J\).  This gives \eqref{inc:q-density}.  The \(R\)
experiment divides \eqref{inc:undivided-C} by \(p_D(Ds_1)^\ell\) and
attaches the unchanged \(\gamma_U\) kernel for \(A\) and uniform \(J\),
giving \eqref{inc:r-density}.  The divisions are legitimate on the
full-measure set just proved.  Since \(r_*>0\) and is finite there,
\(Q\ll R\).
\end{proof}

\begin{lemma}[The induced Stiefel volume]\label{inc:stiefel}
For \(1\le n\le h\), with the induced Euclidean matrix metric,
\begin{equation}\label{inc:stiefel-volume}
 |\mathcal V_{n,h}|=
 2^{n(n-1)/4}\prod_{j=1}^n|\mathbb S^{h-j}|.
\end{equation}
Consequently, for \eqref{inc:dimensions},
\begin{equation}\label{inc:kappa}
 \kappa=\pi^{-n\ell/2}\prod_{j=1}^n
 \frac{\Gamma((p-j+1)/2)}{\Gamma((p-j+1-\ell)/2)}
 \ge\left(\frac{p-n+1-\ell}{2\pi}\right)^{n\ell/2}.
\end{equation}
\end{lemma}

\begin{proof}
For \(n=1\), the frame manifold is \(\mathbb S^{h-1}\), giving the base
case of \eqref{inc:stiefel-volume}.  For \(n\ge2\), the map from
\(\mathcal V_{n,h}\) to its first row in
\(\mathbb S^{h-1}\) has constant normal Jacobian
\(2^{-(n-1)/2}\).  To compute it, move the frame by an orthogonal
transformation to the first \(n\) coordinate rows.  Tangent matrices
then consist of a skew-symmetric block in the first \(n\) columns and
an arbitrary remaining block.  Each of the \(n-1\) skew coordinates
affecting the first row appears twice in the squared matrix norm, so
the corresponding singular value of the first-row map is \(2^{-1/2}\).
The other first-row coordinates have singular value one.  The fiber
is isometric to \(\mathcal V_{n-1,h-1}\).
Slicing thus gives
\[
 |\mathcal V_{n,h}|=
 2^{(n-1)/2}|\mathbb S^{h-1}|\,|\mathcal V_{n-1,h-1}|.
\]
Iteration proves \eqref{inc:stiefel-volume}.  The metric factor depends
only on \(n\), so it cancels in the ratio defining \(\kappa\).
The sphere-area formula
\(|\mathbb S^{h-1}|=2\pi^{h/2}/\Gamma(h/2)\) gives the product in
\eqref{inc:kappa}.

For \(a>0\) and \(\ell/2\ge1\), a gamma random variable \(X\) of shape
\(a\) and scale one satisfies
\(\Gamma(a+\ell/2)/\Gamma(a)=\mathbb E X^{\ell/2}
\ge(\mathbb E X)^{\ell/2}=a^{\ell/2}\) by Jensen's inequality.
Apply this with \(a=(p-j+1-\ell)/2\) and then use
\(a\ge(p-n+1-\ell)/2\).
\end{proof}

\subsection{An inverse kernel at two scales}
\label{inc:inverse-section}

The density ratio in Lemma~\ref{inc:densities} involves
\(p_D(Ds_1)/p_A(As_1)\).  It has no useful pointwise bound.  We control
its positive logarithm instead, using the fact that \(A\) annihilates
the same difference space that governs the conditional Gaussian law
of \(D\).  For \(u>0\), write \(\log^+u=\max\{0,\log u\}\), and set
\(\log^+0=0\).

\begin{lemma}[Two-scale inverse kernel]\label{inc:kernel}
Under \(Q\), define
\[
 \begin{split}
 K&=\int
   \operatorname{dist}(z-s_1,\operatorname{col}U)^{-n}\,d\mu(z),\\
 p_A^*(t)&=\sup_{\rho>0}
   \frac{\mu\{z:\|Az-t\|<\rho\}}{v_n\rho^n},\qquad
 K_0=v_n\|A\|_{\mathrm{op}}^n p_A^*(As_1).
 \end{split}
\]
Then \(K_0\ge2^{-n}\), and
\begin{equation}\label{inc:kernel-bound}
 K\le K_0L_F,\qquad
 L_F=1+\frac{n}{d-\ell-n}+n\log(2/\tau),\qquad
 \tau=\left(\frac{2^{-n}}{Fd^d}\right)^{1/(d-\ell-n)}.
\end{equation}
In particular,
\(\log L_F\le C\log d+C\log(1+\log F)\), and \(K<\infty\).
Moreover,
\begin{equation}\label{inc:conditional-density}
 \mathbb E_Q[p_D(Ds_1)\mid\mathbf s,A,y,J]
 \le (2\pi)^{-n/2}K .
\end{equation}
\end{lemma}

\begin{proof}
For fixed \(z\), \eqref{inc:conditional-D} makes the coordinates of
\(D(z-s_1)\) independent centered Gaussians with variance
\(\operatorname{dist}(z-s_1,\operatorname{col}U)^2\).  Hence
\[
 \frac{\Pr_Q\{\|D(z-s_1)\|<\rho\mid\mathbf s,A,y,J\}}{v_n\rho^n}
 \le (2\pi)^{-n/2}
       \operatorname{dist}(z-s_1,\operatorname{col}U)^{-n}.
\]
The right side is interpreted as \(+\infty\) when the distance is zero.
Integrate in \(z\) using Tonelli's theorem and then apply conditional
Fatou along \(\rho=2^{-j}\).  Lemma~\ref{inc:densities} and
\(Q\ll\mathsf H\) identify the almost-sure limit of the ball averages
with \(p_D(Ds_1)\).  This proves \eqref{inc:conditional-density}
without evaluating an arbitrary density version at a dependent point.

Let \(H(t)\) be the \(\mu\)-probability that the distance in \(K\) is
at most \(t\).  Since \(AU=0\),
\[
 \operatorname{dist}(z-s_1,\operatorname{col}U)\le t
 \quad\Longrightarrow\quad
 \|A(z-s_1)\|\le\|A\|_{\mathrm{op}}t.
\]
The projection law has a Lebesgue density, so projection-ball
boundaries have zero mass.  The definition of \(p_A^*\) therefore gives
\(H(t)\le K_0t^n\).  All these distances are at most two; taking
\(t\downarrow2\) from above also gives \(K_0\ge2^{-n}\).

A second estimate uses the projection onto \(U^\perp\), of dimension
\(d-\ell\le d-2\).  Applying Lemma~\ref{inc:projection} to a matrix
with orthonormal rows for this projection bounds its density by
\(Fd^d/v_{d-\ell}\).  A ball of radius \(t\) centered at the projection
of \(s_1\) consequently has probability at most
\(Fd^dt^{d-\ell}\).  Thus
\begin{equation}\label{inc:two-scales}
 H(t)\le\min\{K_0t^n,\;Fd^dt^{d-\ell}\}.
\end{equation}
The exponent \(d-\ell\) uses the bounded projection density and the
condition \(\ell\ge2\).

The layer-cake formula, including the possibility of an initially
infinite integral, is
\[
 K=2^{-n}+n\int_0^2H(t)t^{-n-1}\,dt.
\]
Here \(d-\ell-n>0\) and \(\tau\le1\).  Use the second estimate in
\eqref{inc:two-scales} up to \(\tau\) and the first thereafter:
\[
 \begin{split}
 K&\le2^{-n}
   +\frac{nFd^d\tau^{d-\ell-n}}{d-\ell-n}
   +nK_0\log(2/\tau)\\
 &\le K_0\left(1+\frac{n}{d-\ell-n}+n\log(2/\tau)\right).
 \end{split}
\]
The last inequality uses
\(Fd^d\tau^{d-\ell-n}=2^{-n}\le K_0\).  Finally,
\[
 \log(1/\tau)
 =\frac{n\log2+\log F+d\log d}{d-\ell-n}.
\]
The dimension ratios in \eqref{inc:dimensions} give
\(L_F\le C(1+d\log d+\log F)\).  Taking its logarithm gives the
stated bound and finiteness of the inverse-distance integral \(K\).
\end{proof}

\begin{lemma}[Maximal average at an actual label]\label{inc:maximal}
For every full-rank \(A\),
\begin{equation}\label{inc:maximal-bound}
 \int\log\frac{p_A^*(As)}{p_A(As)}\,d\mu(s)
 \le Cd+C\log(1+\log F).
\end{equation}
The integrand is nonnegative almost surely.
\end{lemma}

\begin{proof}
Lebesgue differentiation gives \(p_A^*(t)\ge p_A(t)\) for
\(A_\#\mu\)-almost every label.  We prove the particular
centered maximal inequality we need.  Its one-dimensional ancestry is
the maximal-function theory of Hardy and
Littlewood~\cite[Section~III]{HardyLittlewood1930}; for the dimension
cost of customary covering arguments, see Stein and
Str\"omberg~\cite[p.~259]{SteinStromberg1983}.

For \(a>0\), a compact subset of \(\{t:p_A^*(t)>a\}\) has a finite
cover by balls centered in that set, each of projection mass greater
than \(a\) times its volume.  Select disjoint balls in decreasing order
of radius.  Every discarded ball is contained in the triple of a
selected ball, since it meets a selected ball of at least its radius.
The selected masses sum to at most one, so the compact set has volume
at most \(3^n/a\).  Inner regularity gives
\begin{equation}\label{inc:weak-maximal}
 |\{t:p_A^*(t)>a\}|\le3^n/a\qquad(a>0).
\end{equation}
The function is measurable: the supremum may be restricted to positive
rational radii, because a density gives zero mass to ball boundaries.

Let \(E_A\) be the projection ellipsoid and \(V_A=|E_A|\).
Since \(p_A^*\le\|p_A\|_\infty\), integration of level sets, using
the bound \(V_A\) up to level \(1/V_A\) and
\eqref{inc:weak-maximal} above that level, gives
\[
 \int_{E_A}p_A^*(t)\,dt
 \le1+3^n\log^+(V_A\|p_A\|_\infty)
 \le1+3^n(\log F+d\log d).
\]
Jensen's inequality with respect to the probability density \(p_A\)
then yields
\[
 \int p_A(t)\log\frac{p_A^*(t)}{p_A(t)}\,dt
 \le\log\int_{\{p_A>0\}}p_A^*(t)\,dt
 \le\log\int_{E_A}p_A^*(t)\,dt.
\]
Values on \(\{p_A=0\}\) do not enter the integral.  The preceding
estimate and \(n\asymp d\) prove \eqref{inc:maximal-bound}.
\end{proof}

\begin{proof}[Proof of Proposition~\ref{inc:block}]
On the full-measure set in Lemma~\ref{inc:densities}, the log likelihood
ratio is the sum of the message term \(\log(w\phi_J(A,As_1))\), the term
\begin{equation}\label{inc:gram-term}
 -\frac{n\ell}{2}\log(2\pi)-\log\kappa+
 \frac n2\log
 \frac{\det(U^{\mathsf T}B^{\mathsf T}BU)}{\det(U^{\mathsf T}U)},
\end{equation}
and \(\ell\log(p_D(Ds_1)/p_A(As_1))\).
The message term is at most \(\log w\).

Under \(Q\), \(B\) is independent of \(U\).  Conditional on \(U\), put
\(V=U(U^{\mathsf T}U)^{-1/2}\).  Then \(G=BV\) is a standard Gaussian
\(p\times\ell\) matrix and
\[
 \frac{\det(U^{\mathsf T}B^{\mathsf T}BU)}{\det(U^{\mathsf T}U)}
 =\det(G^{\mathsf T}G).
\]
The lower bound in \eqref{inc:kappa} and the determinant--trace
inequality bound \eqref{inc:gram-term} above by
\[
 \frac{n\ell}{2}\log
 \frac{\operatorname{tr}(G^{\mathsf T}G)}
      {\ell(p-n+1-\ell)}.
\]
Its positive part has expectation \(O(n\ell)\).  Indeed,
\(\mathbb E\operatorname{tr}(G^{\mathsf T}G)=p\ell\), the ratio
\(p/(p-n+1-\ell)\) is bounded, and
\(\mathbb E\log^+X\le\log(1+\mathbb E X)\) for \(X\ge0\).

For the remaining ratio, use
\(\log^+u\le\log(1+u)\), conditional Jensen, and
\eqref{inc:conditional-density} and \eqref{inc:kernel-bound}.  Since
\(p_A^*(As_1)/p_A(As_1)\ge1\) almost surely, this gives
\begin{align}
 &\mathbb E_Q\!\left[
   \log^+\frac{p_D(Ds_1)}{p_A(As_1)}
   \,\middle|\,\mathbf s,A,y,J\right]\notag\\
 &\quad\le \log2+
  \log^+\!\bigl((2\pi)^{-n/2}v_n\|A\|_{\mathrm{op}}^n\bigr)
  +\log L_F+\log\frac{p_A^*(As_1)}{p_A(As_1)}.
 \label{inc:log-ratio-conditional}
\end{align}
The Gaussian operator term has expectation \(O(n)\).  Here are the
elementary estimates.  A \(1/4\)-net of the unit sphere in
\(\mathbb R^h\) has size at most \(9^h\), by packing disjoint radius
\(1/8\) balls.  Nets in the two dimensions bound \(\|A\|_{\mathrm{op}}\)
by twice the maximum of at most \(9^{n+d}\) absolute standard Gaussian
bilinear forms.  A Gaussian tail bound gives
\(\mathbb E\|A\|_{\mathrm{op}}\le C\sqrt d\).
Also, integrating the Gaussian density over the ball of radius
\(\sqrt n\) gives \(v_n\le(2\pi e/n)^{n/2}\).  Therefore
\[
 \log^+\!\bigl((2\pi)^{-n/2}v_n\|A\|_{\mathrm{op}}^n\bigr)
 \le n\log(1+\sqrt{e/n}\,\|A\|_{\mathrm{op}}),
\]
whose expectation is \(O(n)\) by Jensen and \(n\asymp d\).

The \((A,s_1)\)-marginal under \(Q\) is the original independent
Gaussian/prior draw.  Averaging \eqref{inc:log-ratio-conditional} and
using Lemmas~\ref{inc:kernel} and~\ref{inc:maximal} consequently gives
\begin{equation}\label{inc:log-ratio}
 \mathbb E_Q\log^+\frac{p_D(Ds_1)}{p_A(As_1)}
 \le Cd+C\log(1+\log F).
\end{equation}
All positive parts in the log likelihood ratio are now integrable.
Since \(Q\ll R\), its negative part is integrable as well: with
\(h=dQ/dR\),
\(\int_{\{h<1\}}-h\log h\,dR\le1/e\).
Thus no undefined subtraction occurs, and
\[
 \operatorname{KL}(Q\|R)
 \le\log w+Ckd+Ck\log(1+\log F).
\]
Combine this with \eqref{inc:entropy-lower} and divide by \(k\).
\end{proof}

\subsection{Finite histories and the precision endpoint}
\label{inc:history-section}

We state the consequence first for average success under the uniform
spherical prior, using the exact finite-state randomized model in
Definition~\ref{def:streaming-model}.

\begin{theorem}[Average-prior precision bound]\label{inc:precision}
There is a universal \(c>0\) with the following property.  Let
\(M=M(d)\) be a nonnegative integer with \(M(d)=o(d^2)\).
For all sufficiently large \(d\), with the
threshold allowed to depend on the rate \(M(d)/d^2\to0\), suppose
\(S\sim\sigma\) and a learner in Definition~\ref{def:streaming-model} stops by a
deterministic integer horizon \(T\ge0\).  If
\[
 \Pr\{\angle(\widehat S,S)\le\epsilon\}\ge2/3,
 \qquad 0<\epsilon\le1/10,
\]
then
\begin{equation}\label{inc:precision-bound}
 T\ge c\,d\log(1/\epsilon).
\end{equation}
The bound permits arbitrary computation within a transition.
\end{theorem}

For the proof of Theorem~\ref{inc:precision} and its auxiliary lemmas
below, apply the fixed-shared-seed part of Lemma~\ref{lem:padding} to the
stated learner in Definition~\ref{def:streaming-model}.  We work in the
resulting Borel experiment with data-independent initialization and
uniform-prior success at least \(2/3\), retaining fresh transition and
output randomness.  This is the standing reduction for this subsection.

\begin{lemma}[Short-horizon reduction]\label{inc:short}
Put \(L_\epsilon=\log(1/\epsilon)\).  Pad stopping to \(T\) by retaining
the stopped state and its stopping index.  At most
\begin{equation}\label{inc:width}
 w=(T+2)2^M
\end{equation}
padded states suffice.  Under the hypotheses of
Theorem~\ref{inc:precision}, if \(T<dL_\epsilon\), then
\begin{equation}\label{inc:short-bounds}
 L_\epsilon=O(1+M/d)=o(d),\qquad
 T=O(d^2),\qquad \log w=o(d^2).
\end{equation}
\end{lemma}

\begin{proof}
The padding count in Lemma~\ref{lem:padding}, including stopping at
index zero, gives \eqref{inc:width}.  Its stopped records remain fixed
while fresh unused observations are generated.
The output-capacity bound in Lemma~\ref{lem:early-clock}, with success
at least \(2/3\) and \((T+1)2^M\le w\), gives
\[
 \frac23\le w\epsilon^{d-1},\qquad
 (d-1)L_\epsilon\le
 M\log2+\log(T+2)+\log(3/2).
\]
Because \(L_\epsilon\ge\log10\) and \(T<dL_\epsilon\),
\(\log(T+2)\le\log d+\log L_\epsilon+C\).
Using \(\log L_\epsilon\le L_\epsilon\) yields
\[
 (d-2)L_\epsilon\le M\log2+\log d+C.
\]
This proves \(L_\epsilon=O(1+M/d)=o(d)\).  The remaining assertions
follow from \(T<dL_\epsilon\) and \(M=o(d^2)\).
\end{proof}

\begin{lemma}[Information accumulated over histories]\label{inc:history}
In the short-horizon case of Lemma~\ref{inc:short}, put
\(b=\lceil T/n\rceil\) and pad with fresh unused samples to \(bn\)
observations.  Let \(V_i\) be the transcript of the first \(i\)
block-end padded states, with \(V_0\) trivial.  Let
\(B\sim\gamma_{p,d}\) be independent of the entire learner experiment
and used only for analysis.  Then, for all sufficiently large \(d\),
\begin{equation}\label{inc:history-bound}
 I(S;V_b\mid B,BS)\le Cbd
\end{equation}
for an absolute \(C\).  The dimension threshold may depend on the
rate \(M/d^2\to0\).
\end{lemma}

\begin{proof}
For a positive-probability history \(V_{i-1}=v\),
Lemma~\ref{lem:conditional-priors} gives the Bayes density
\[
 f_v(s)=\frac{\Pr\{V_{i-1}=v\mid S=s\}}{\Pr\{V_{i-1}=v\}}
 \le F_v:=\frac1{\Pr\{V_{i-1}=v\}}
\]
relative to \(\sigma\).  A Borel version is available from the Borel
learner.  The history fixes the previous padded state.  Its update
over the next fresh block is a kernel of \((A,AS)\) to at most \(w\)
next states.  For the first block, independent initialization can be
integrated into the same kernel.  The matrix \(B\) is absent from the
learner and independent of its entire experiment; it remains independent
of the entire conditional signal, fresh-block, and message experiment
after conditioning on \(v\).
Proposition~\ref{inc:block} therefore applies at each such history.

The finite history has entropy
\(H(V_{i-1})\le b\log w\).  Hence Jensen's inequality gives
\[
 \mathbb E_v\log(1+\log F_v)
 \le\log(1+H(V_{i-1}))
 \le\log(1+b\log w)=O(\log d).
\]
The last bound uses \(b=O(d)\) and \(\log w=O(d^2)\) eventually, from
\eqref{inc:short-bounds} and \(n\asymp d\).  The same bounds give
\(\log w/k=o(d)\).  Averaging the one-block estimate thus bounds each
increment
\(I(S;V_i\mid V_{i-1},B,BS)\) by \(Cd\) eventually.
The conditional chain rule, with the same independent \(B\) throughout,
proves \eqref{inc:history-bound}.
\end{proof}

\begin{lemma}[Precision inside an independent fiber]\label{inc:fiber-precision}
Let \(S\sim\sigma\), and let a unit-vector output \(\widehat S\) have
angular error at most \(\epsilon\le1/10\) with probability at least
\(2/3\).  Let \(B\sim\gamma_{p,d}\), \(p=\lfloor d/4\rfloor\), be
independent of the entire experiment producing \((S,\widehat S)\).
Then, for all sufficiently large \(d\),
\begin{equation}\label{inc:fiber-bound}
 I(S;\widehat S\mid B,BS)
 \ge \frac{d-p-1}{3}\log\frac1{4\epsilon}-\log2
 \ge c_1dL_\epsilon
\end{equation}
with an absolute \(c_1>0\).  In the learner model above, success
probability at least \(2/3\) is also impossible for \(T\le d/4\),
even if all observations are retained.
\end{lemma}

\begin{proof}
Apply Lemma~\ref{lem:residual-information} with \(V=\widehat S\), the
identity output kernel, \(b=p\), and \(r_0=1/2\).  For the residual radius
\(R=\|P_{\ker B}S\|\), Lemma~\ref{lem:residual-sphere} gives
\[
 \Pr\{R<1/2\}\le\frac{4p}{3d}\le\frac13.
\]
Success at least \(2/3\) therefore leaves event mass at least \(1/3\),
and \eqref{eq:residual-information} gives the first inequality in
\eqref{inc:fiber-bound}.  The second follows from
\[
 \log\frac1{4\epsilon}
 \ge\left(1-\frac{\log4}{\log10}\right)L_\epsilon,
\]
because \(d-p-1\) is an absolute positive fraction of \(d\); the
subtractive \(\log2\) is absorbed for large \(d\).

For the final assertion, grant the learner all pre-generated observations
through its horizon, including ignored rows after stopping.  Its output
is a kernel of those data and independent randomness.  Lemma~\ref{lem:raw-sample-bound}
shows that success at least \(3/5\), and hence at least \(2/3\), already
requires \(T>\lfloor d/2\rfloor\) for large \(d\).  In particular,
\(T>d/4\).  That lemma includes the deterministic zero-row residual
sphere, so the assertion covers \(T=0\) as well.
\end{proof}

\begin{proof}[Proof of Theorem~\ref{inc:precision}]
If \(T\ge dL_\epsilon\), the desired bound is immediate.  Otherwise
apply Lemmas~\ref{inc:short} and~\ref{inc:history}.  The terminal output
is a kernel of the final padded state, hence of \(V_b\); conditional
data processing and Lemma~\ref{inc:fiber-precision} give
\[
 Cbd\ge I(S;V_b\mid B,BS)
 \ge I(S;\widehat S\mid B,BS)
 \ge c_1dL_\epsilon.
\]
Thus \(b\ge c_0L_\epsilon\) for an absolute \(c_0>0\).
If \(L_\epsilon\ge2/c_0\), then
\[
 T>n(b-1)\ge(c_0/2)nL_\epsilon\ge c\,dL_\epsilon
\]
after adjusting an absolute constant.  If \(L_\epsilon<2/c_0\), the
separate conclusion \(T>d/4\) of Lemma~\ref{inc:fiber-precision}
gives the same order after another absolute adjustment.  This proves
\eqref{inc:precision-bound} uniformly over the allowed \(\epsilon\).
\end{proof}

A guarantee of success probability at least \(2/3\) for every fixed
unit signal implies the uniform-prior average guarantee, so it has the
same consequence.  The argument uses fresh exact labels, a finite
persistent state alphabet, and a deterministic horizon; its additional
projections and tuple resampling belong only to the analysis.
